The chronological compression after the close of intensive reported sampling is decisive.

The Year-5 reporting period ends 31 May 2019. From that date until the first detectable human cases in Wuhan in late 2019 lies a window of roughly six months. Within that narrow interval any pure natural-proximal account must either produce the missing FCS-bearing specimen or explain its absence while still generating a fully human-adapted virus that appears, fully formed, at a single geographic point.

No such specimen was reported. No bat isolate, no intermediate-host isolate, and no environmental sample recovered from the surveyed southern reservoirs during or immediately after that window has ever been shown to carry the exact PRRA insertion with the CGG-CGG codon pair. Subsequent sampling has continued to recover nearer relatives that still lack the configuration. The molecular clock ensures that later material cannot be retrofitted into a pristine 2019 proximal ancestor; the evolutionary time that has already elapsed is irreversible.

The alternative—positing a rapid cascade of multi-species jumps, hybridizations, or chimeric recombinations sufficient to assemble the defining genomic features inside six months—requires an extraordinary sequence of low-probability events: successive cross-species transmissions among multiple hosts, repeated recombination, fixation of the precise furin-cleavage insertion, and silent amplification, all occurring outside the intensive surveillance network and leaving no recoverable intermediate. That cascade multiplies unobserved entities and compresses evolutionary processes into a temporal span that is itself an outlier relative to the documented natural history of related sarbecoviruses.

No compensatory mechanism exists. One cannot later “make up” a fake bat or invent a temporally matching natural progenitor once the molecular clock has advanced. The absence therefore remains non-neutral: it permanently elevates the evidential burden on every pure natural-proximal narrative that depends on an unsampled FCS-bearing form circulating in the studied reservoirs.

By contrast, the research-related pathway faces no such requirement. Samples had already moved from high-risk field sites into a laboratory continuum engaged in sequencing, culture, and experimental work on SARS-related coronaviruses. The geographic terminus of that continuum was Wuhan; the temporal terminus of intensive reported sampling was mid-2019. An introduction arising inside that continuum aligns with the single-point origin signature without needing a rapid multi-species cascade or a missing natural specimen. The pathway remains unproven, yet it is not constrained by the same chronological and molecular-clock barriers that confront pure natural accounts.

In short, the six-month window after May 2019 does not merely leave a gap; it closes the practical possibility of recovering a temporally authentic natural proximal ancestor while forcing pure natural explanations into ever more elaborate unobserved sequences. The research-related alternative stays coherent precisely because it does not require those sequences.
